\section{Comparison between PLONK and our SNARKs}
\label{suplementary_plonk_comparison}

\noindent In the following, we briefly look at the differences between PLONK universal SNARK and the SNARKs designed in this work. We observe 
that while the NP relation that defines PLONK is more general, the relations that define our SNARKs are bespoke as we are only interested in efficiently proving 
public key aggregation. Because our relations are so bespoke, it turns out we do not require the full functionality that PLONK has to offer, and, in particular, our SNARKs  
do not require any permutation argument. \\

\noindent A second difference is that while PLONK's circuit is defined by a number of selector 
polynomials (which are a type of pre-processed polynomials) and a PLONK verifier needs to perform 
a one-time expensive computation of the polynomial commitments to those selector polynomials, our SNARK verifiers 
are able to avoid such a pre-processing phase. Indeed, in the case of $\mathscr{P}^h_{\mathsf{a}}$ (which is the only one of 
our two SNARKs that has a polynomial, namely $\mathit{aux}(X)$, that defines its circuit), our respective SNARK verifier does not need to compute a 
commitment to its only ``selector polynomial'' as, due to its structure, $\mathit{aux}(X)$ can be directly and efficiently evaluated by our SNARK verifier itself. \\ 

\noindent A third difference is that using our two-steps compiler, our SNARKs verifiers are able to efficiently handle input vectors of length $O(n)$, 
where the degree of the polynomials committed to by our SNARK provers is also $O(n)$. Our SNARKs verifiers achieve efficiency by offloading the 
expensive polynomial commitment computation involving the public inputs to a trusted third party. \\

\noindent Moreover, while PLONK does not incorporate trusted inputs, one can easily apply the Step 2 of our compiler to PLONK. In particular, one could imagine a situation 
where a PLONK verifier is relying on a trusted party to compute some or all of the polynomial commitments to the circuit's selector polynomials. This is equivalent to our hybrid 
model SNARK definition applied to PLONK. The benefit is that by delegating such a computation, the PLONK verifier becomes more efficient. \\ %at the expense of relying on some trust assumptions. \\

\begin{comment}
The next paragraph should be commented out if we do not include in the submission the light client model instantiation section, i.e., section 5.3 on eprint.
\end{comment}

\noindent Finally, looking at our light client system instantiation in Appendix~\ref{sec:LCinstantiation} due to the inductive structure of the soundness proof (Theorem~\ref{th:soundness}), 
the efficiency of using a hybrid model SNARK has an even greater impact for the light client system verifier than that compared to verifying multiple instances of PLONK for the same circuit:
while for the latter the PLONK verifier has to compute commitments to selector polynomial only once anyway, in the case of the former, the commitments to public inputs may differ 
at every step; hence, a trusted third party relives a higher computation burden from the light client verifier overall. 


